\section{Extensive games, realization and strategic interpretation}
\label{annex:p17-rationality-9-1}

\begin{annexitem}{annex:p17.rat91.tree.solutions}{thm:p17-rat91-tree-solutions}{Every pure Nash cell and backward comparison}
\begin{proof}
In the first normal form row's best replies to \(WY,WZ,XY\)
are each \(AC\); to \(XZ\) they are \(BC,BD\).
Column's best replies to \(AC\) are \(XY,XZ\), to
\(AD\) they are \(WY,WZ\), and to \(BC,BD\) they
are \(WZ,XZ\). Intersection gives exactly
\((AC,XY),(BC,XZ),(BD,XZ)\).
Backward maximization selects \(C\) over \(D\) by \(5>4\),
\(Z\) over \(Y\) by \(4>3\), \(X\) over \(W\)
after \(C\) by \(1>0\), then \(B\) over \(A\)
by \(3>2\). These strict comparisons force the unique
behavioral subgame-perfect profile \((BC,XZ)\).
The other two pure equilibria fail at off-path \(Y\)
and \(D\), respectively.

In the second normal form row's best replies to \(W,X\)
are \(AC\) and \(BC,BD\), respectively. Column's
best replies to \(AC,AD,BC,BD\) are \(X,W,X,X\),
giving exactly \((BC,X),(BD,X)\). The last player-1
subgame forces \(C\). The reduced \(A,B\) versus
\(W,X\) matrix has cells \((5,0),(2,1),(0,3),(3,4)\).
Column \(X\) dominates \(W\) by 1 against every row
mixture; against \(X\), row strictly prefers \(B\).
Thus the full behavioral subgame-perfect profile is uniquely
\((BC,X)\), not merely unique on the pure list.
\end{proof}
\end{annexitem}

\begin{annexitem}{annex:p17.rat91.outside}{thm:p17-rat91-outside}{Outside option and the explicit forward rule}
\begin{proof}
Row's best replies to \(L,R\) are respectively \(SU\)
and \(XU,XD\). Column's best replies to
\(SU,SD,XU,XD\) are \(L,R,\{L,R\},\{L,R\}\).
This gives the three stated pure Nash cells.
The proper subgame must have one of the three complete
Battle-of-the-Sexes equilibria. Its row payoffs are
4, 2 and \(4/3\), compared with exit's 3. The strict
root comparisons force \((SU,L),(XD,R),(XM,N)\),
with no additional initial mixtures. Their full payoffs
are \((4,2),(3,3),(3,3)\). In \((XU,R)\) the
subgame continuation \((U,R)\) is not Nash: changing
to \(L\) raises column from 0 to 2, and changing
to \(D\) raises row from 0 to 2.

Plan \(SD\) pays 0 against \(L\), 2 against \(R\).
Plan \(SM\) pays \(8/3\) against \(L\), \(2/3\)
against \(R\). An exit plan always pays 3 and strictly
dominates both; linearity gives the same strict comparisons
against every column mixture. Under the stipulated explanatory
exclusion rule, observed \(S\) therefore cannot be attributed
to either plan. The \((XD,R)\) candidate would prescribe
\(SD\) after unexpected start, while \((XM,N)\) would
prescribe \(SM\). Only \((SU,L)\) remains among the
subgame-perfect candidates. This proves that particular
forward contrast, not a general forward-induction theorem.
\end{proof}
\end{annexitem}

\begin{annexitem}{annex:p17.rat91.horse}{thm:p17-rat91-horse}{All unilateral tests and the off-path belief threshold}
\begin{proof}
At \((D,R,R)\), payoff \((3,2,2)\), player 1's
\(R\) deviation gets 1 instead of 3; player 2's
node remains unreached under either own instruction;
player 3's \(L\) gets 0 instead of 2. Thus no pure
or mixed unilateral deviation gains. At player 2's
hypothetical node, however, \(R\) gives 1 while
\(D\) followed by player 3's \(R\) gives 4.
A subtree there cuts player 3's shared information set,
and the shared set is not itself a singleton root.
There is no proper subgame, so this Nash profile is
subgame perfect despite that local failure.

At \((R,R,L)\), payoff \((1,1,1)\), player 1's
\(D\) gets 0 and player 2's \(D\) gets 0.
Player 3 is unreached and either own instruction leaves
its payoff 1. This too is Nash and subgame perfect.
At its shared set, belief \(\mu\) on the left node
gives player 3 \(1-\mu\) from \(L\), \(2\mu\)
from \(R\). The difference \(3\mu-1\) proves the
threshold, with both actions tied at \(\mu=1/3\).
The whole set has equilibrium probability zero, so
ordinary conditioning cannot define \(\mu\).
\end{proof}
\end{annexitem}

\begin{annexitem}{annex:p17.rat91.memory}{thm:p17-rat91-memory}{Memory calculations and behavioral nonexistence}
\begin{proof}
Against \(XX\) in Figure 9.1.6a, either initial
move ends at \((1,1)\), and player 1's last
instruction is unreached. No whole-plan deviation
improves its payoff. Against \(LL\), player 2's
reached left node gives 1 from \(X\), 0 from
\(D\); its right node is unreached. Thus
\((LL,XX)\) is Nash. If first and last roles
are separate players both receiving player 1's
utility, the first gets 1 from either move and
the last is never reached; \((L,XX,L)\) is
Nash in that distinct model too.
The asserted \((LL,DX)\) instead follows
\(LDL\), paying \((3,0)\). Changing only
player 2's left instruction to \(X\) follows
\(LX\), paying \((1,1)\); player 2 strictly
gains. This also refutes the separate-role version.

With player 2 fixed at \(W\) in Figure 9.1.6b,
fresh fair \(L/R\) draws give probability
\(1/2\) to root \(R\) and \(1/4\) each
to \(LWL,LWR\). The expectation is
\[
 \tfrac12(3,4)+\tfrac14(5,0)+\tfrac14(0,2)
                              =(11/4,5/2).
\]
The fair mixture of always-\(L\) and always-\(R\)
gives instead \(\tfrac12(5,0)+\tfrac12(3,4)=(4,2)\).
Its complete-plan draw correlates repeated instructions,
unlike independent fresh behavioral draws.

For the existence countercase, player 2 first chooses
\(W\) or \(R\), unobserved by player 1. Player 1
then makes two successive \(L/R\) choices regardless
of its first choice, at one information set containing
all its first- and second-choice nodes. Every play has
exactly three moves. If its two choices are \(LL\),
player 1 gets 1 after \(W\), \(-1\) after player
2's \(R\); for \(RR\) these values are \(-1,1\).
If the choices differ, player 1 gets \(-2\) after
either opponent move. Player 2 gets the negative
at every terminal node. This specifies the entire
finite tree and payoffs; its repeated information
set violates perfect recall.

Player 1's two pure plans are always-\(L\) and
always-\(R\), inducing \(LL,RR\); their mixed
normal form is Matching Pennies, with fair mixed
Nash equilibrium. A behavioral strategy draws
\(L\) with one fixed probability \(p\) at each
visit, freshly. Let \(q\) be player 2's probability
of \(W\). Then player 1's payoff is
\[
 V(p,q)=(2q-1)[p^2-(1-p)^2]-4p(1-p)
       =(2q-1)(2p-1)-4p(1-p).
\]
For \(0<p<1\) the first term interpolates endpoint
payoffs and the second is strictly negative. Hence
\(V(p,q)<\max\{V(0,q),V(1,q)\}\), forcing
every behavioral best reply of player 1 to be an
endpoint. At \(p=0\), player 2 maximizes
\(-V=2q-1\) by \(q=1\), but then player 1
strictly prefers \(p=1\). At \(p=1\), player
2 maximizes \(-V=1-2q\) by \(q=0\), but
then player 1 strictly prefers \(p=0\). Both
endpoint cases fail. No behavioral Nash exists,
and therefore no behavioral subgame-perfect profile
exists. This explicit best-response contradiction,
not a simulation, proves the finiteness boundary.
\end{proof}
\end{annexitem}

\begin{annexitem}{annex:p17.rat91.blackmail}{thm:p17-rat91-blackmail}{Credible continuation and the explicit commitment move}
\begin{proof}
The original matrix is
\[
\begin{array}{c|cc}&R&P\\\hline
 X&(0,0)&(3,1)\\S&(2,2)&(2,2).
\end{array}
\]
Eve's best replies to \(R,P\) are \(S,X\).
Adam's best reply to \(X\) is \(P\), and
both replies tie against \(S\), giving exactly
the two pure Nash cells. The subgame after
\(X\) strictly requires \(P\), since \(1>0\).
Eve then strictly prefers \(X\), since \(3>2\).
Thus \((X,P)\) is the unique behavioral
subgame-perfect profile.
In the extended game this continuation after
\(N\) gives Adam 1. After \(C\), Eve chooses
\(S\), since \(2>0\), giving Adam 2. Adam
initially chooses \(C\), since \(2>1\).
All comparisons are strict, proving uniqueness
of the complete profile \((CP,XS)\).
Without the added commitment node, the original
report/pay comparison \(0<1\) is unchanged;
an unmodeled declaration cannot supply the
new game's credibility.
\end{proof}
\end{annexitem}

\begin{annexitem}{annex:p17.rat91.repetition}{thm:p17-rat91-repetition}{Finite backward reasoning and the exact trigger condition}
\begin{proof}
In the last finite stage, \(D\) beats \(U\)
and \(R\) beats \(L\) by 1 against either
opponent action. No future payoff remains, so
every subgame-perfect profile defects at every
last-stage history. In the previous stage its
future payoff is now the constant \((1,1)\)
whatever current choices occur; adding this
constant preserves strict dominance. Backward
induction forces the same choice at every
history of every stage, excluding mixtures.
Conversely all-defection is subgame perfect:
each changed action loses 1 and cannot change
the fixed continuation. At two stages the first
loss is immediate and the remaining payoff
stays 1.

For infinite play bounded payoffs make the
discounted series absolutely convergent. After
any defection the opponent defects forever;
own defection gives 1, cooperation 0, at
every stage. No arbitrary future plan beats
the punishment value 1. At a cooperative
history staying forever gives normalized value
3. A first immediate defection yields 4 now,
then at most 1 at every later stage, so
its largest possible payoff is
\[
 (1-\delta)4+(1-\delta)\sum_{t\geq1}\delta^t
                             =4-3\delta.
\]
If first defection occurs at stage \(k\),
the largest possible payoff is
\(3(1-\delta^k)+\delta^k(4-3\delta)\).
For \(\delta\geq1/3\) this is at most 3
for every finite \(k\); never defecting
gives 3. Randomized deviations average over
their first-defection time, including infinite
time, and preserve this bound. Thus all
unilateral strategies, not just one-stage
deviations, are excluded at every cooperative
history. Combined with the punishment argument
this proves subgame perfection at all finite
histories. If \(\delta<1/3\), immediate
defection gives \(4-3\delta>3\) and Nash
fails. Equality permits but does not force
cooperation. This establishes only the
specified trigger profile under its stated
monitoring and discount assumptions.
\end{proof}
\end{annexitem}

\begin{annexitem}{annex:p17.rat91.calculations}{thm:p17-rat91-calculations}{All simultaneous calculations}
\begin{proof}
Independence gives the four cell probabilities
\(pq,p(1-q),(1-p)q,(1-p)(1-q)\), proving the expectation formula.
The differences of first minus second pure actions are
\[
 d_1(q)=q(x_1-x_3)+(1-q)(x_2-x_4),\quad
 d_2(p)=p(y_1-y_2)+(1-p)(y_3-y_4).
\]
A positive difference forces the first action, a negative difference
the second; a zero difference permits any mixture. For Prisoner's
Dilemma both differences are \(-1\), forcing \(p=q=0\),
payoff \((1,1)\), and strictness. For Matching Pennies they are
\(4q-2,2-4p\). A pure choice makes the opponent's unique response
favor the opposite own choice, excluding any equilibrium with one
player pure. Both differences therefore vanish, giving fair mixtures
and payoff zero.

For Battle of the Sexes the differences are \(6q-2,6p-4\).
The matching corners \((1,1),(0,0)\) are strict equilibria.
If either player is pure the response forces the matching corner;
otherwise both differences vanish, giving \((2/3,1/3)\).
Its \(UL,DR\) cell probabilities are both \(2/9\), so both
payoffs are \(4/3\). For the unnamed game the differences are
\(q-1,2p-1\). If \(q<1\), row must set \(p=0\),
then column must set \(q=0\). If \(q=1\), row may use any
mixture, but column \(L\) is optimal exactly when \(p\geq1/2\).
The payoffs on this segment are \((4,4p)\), while the strict
corner \((0,0)\) gives \((1,1)\). The corner \((1,1)\)
is nonstrict because both row actions give 4. These response
conditions exhaust the domains, establishing completeness.

Changing only Matching Pennies' row \(UL\) payoff to 2 gives
differences \(5q-2,2-4p\). The same boundary-response contradiction
leaves only \((1/2,2/5)\). Both row actions give \(1/5\) and
both column actions give zero. On any maintained nondegenerate
interior branch, the row-coefficient equation determines \(q\)
and the column-coefficient equation determines \(p\), proving
the local sensitivity claim. It makes no global statement across
ties or branch changes. By linearity, a mixed best reply assigning
positive weight to two pure actions makes both pure actions best
replies, so the mixture cannot be the unique best reply.
\end{proof}
\end{annexitem}

\begin{annexitem}{annex:p17.rat91.realization}{thm:p17-rat91-realization}{Full realization bridge, including zero own-history weights}
\begin{proof}
At a node of player \(i\), list its prior own information sets
and moves as \((H_1,a_1),\ldots,(H_k,a_k)\). Perfect recall
makes this the same list throughout a current information set
\(H\), with all sets distinct and \(H\) absent from the list.
The mixed mass of plans prescribing this list is \(r_i(H)\);
add the condition \(s_i(H)=a\) to obtain \(r_i(H,a)\).
The latter events partition the former, so their weights sum
to \(r_i(H)\). The positive-denominator ratios form a probability
law. A zero denominator has all numerators zero and admits any
behavioral law.

Fix a terminal history and list player \(i\)'s own moves on it
as \((H_1,a_1),\ldots,(H_m,a_m)\). Let \(R_k\) be
the mixed mass of plans prescribing the first \(k\) such moves,
with \(R_0=1\). Perfect recall gives
\(r_i(H_k)=R_{k-1}\), \(r_i(H_k,a_k)=R_k\).
If all denominators are positive, the product of behavioral
probabilities telescopes to \(\prod_kR_k/R_{k-1}=R_m\).
If a weight first vanishes at \(R_k\), the previous weight
is positive and the corresponding behavioral factor is zero.
All later weights are zero; the product is again \(R_m=0\),
regardless of arbitrary laws assigned after the first zero.
An empty list gives factor 1 in both descriptions.

The probability of a terminal history factors into Nature's
probabilities and one own-prescription factor per player. For
independent mixtures each factor is the corresponding mixed
mass; for behavioral opponents it is a product of local
probabilities. The identity just proved preserves the replaced
player's factor against either opponent domain, hence preserves
every terminal probability and all expected payoffs.

Conversely assign each pure plan the mass
\(\sigma_i(s_i)=\prod_H b_i(s_i(H)\mid H)\).
Finite product expansion shows total mass one. Summing over
unvisited sets contributes factors 1, leaving the local probability
product along any path. Perfect recall prevents a path from
asking a single information set for two different prescriptions.
Thus terminal laws are again equal against every fixed opponent
profile. Applying either construction to any unilateral deviation
proves equality of its attainable payoffs. This does not determine
a conditional node belief at a null information set.
\end{proof}
\end{annexitem}

\begin{annexitem}{annex:p17.rat91.spe.existence}{thm:p17-rat91-spe-existence}{Finite recursive SPE construction and its deviation proof}
\begin{proof}
Subgames are disjoint or nested: incomparable roots have disjoint
descendant trees, while the later of two comparable roots has its
whole descendant tree inside the earlier one's. A proper subgame
has fewer nodes. Induct on the number of decision nodes, starting
with the trivial Nature-only or terminal case. Solve all maximal
proper subgames recursively and replace each by its saved equilibrium
expected-payoff vector. They are disjoint, so the saved behavioral
profiles are compatible. The residual game retains complete
information sets and unchanged histories at retained nodes, hence
perfect recall.

Its normal form has finite nonempty pure-plan sets, including the
empty plan for a passive player. \cref{thm:p17-nash-1950} supplies
a mixed Nash profile. Behavioralize using the preceding full bridge,
which preserves every unilateral-deviation payoff, and restore
the saved subgame profiles.

For any behavioral deviation, its payoff upon entry to a solved
subgame is at most that subgame's equilibrium value, by the subgame
Nash property. Entry is at one specified node, and opponents'
behavioral draws inside are independent of earlier draws. Changing
one's outside instructions may change entry probabilities but not
this conditional upper bound. Replacing the deviating continuations
by saved continuations raises payoff by a sum of nonnegative
conditional gains weighted by nonnegative entry probabilities;
zero-probability entries contribute zero. Thus the deviating full
payoff is at most the payoff of its induced residual-game deviation,
which residual Nash bounds by the equilibrium payoff. The realization
bridge also excludes arbitrary mixed-plan gains against the restored
behavioral opponents. Every proper subgame is within a maximal
proper one and is Nash by induction; the full game is Nash by the
preceding argument. This proves subgame perfection everywhere.
Strict decrease in node count proves termination.

With singleton information sets, maximize at the last decision nodes
then replace solved continuations by their payoffs and work backward.
Induction on remaining depth shows no unilateral change can improve
in any rooted subgame: later instructions already maximize continuation
and earlier alternatives were included in the maximization. This
gives a pure subgame-perfect profile. Conversely any subgame-perfect
profile must optimize at last nodes and recursively at each earlier
node. Unique local maximizers therefore force the same pure profile
throughout; ties force only support on maximizing moves and do
not imply uniqueness.
\end{proof}
\end{annexitem}
